\section{Introduction}
\label{sec:introduction}

This report is written for the summer semester part of the Simulation Sciences Seminar at RWTH Aachen University, and it was supervised by Anna Ranno. It presents stentFIT, which is a semi-automated pipeline for virtual stent implantation with mixed-dimensional modelling. The pipeline starts with the surface geometry as an input file of a stent and turns it into a 1D beam model which can be used by a solver. Then based on the load cases in the finite element solver, it sets up the simulation and analyzes the mechanical response. At the end of this project, a python packaged called stentFIT is developed and its source is publicly available at \url{https://github.com/VuralAktas/stentFIT} while the documentation of the implementation is at \url{https://stentfit.readthedocs.io}.


\subsection{Motivation}
\label{sec:motivation}

The cardiovascular diseases are one of the leading causes of death worldwide, with millions of people affected each year. According to the Eurostat, in 2023, more than 30\% of deaths in the European Union were caused by cardiovascular diseases, while there were no other diseases that have higher percentage of deaths \citep{eurostat2023}. The most common cardiovascular disease is coronary artery disease, which is caused by the narrowing of the coronary arteries due to plaque buildup, and it can lead to serious health problems such as a heart attack and even death. The most common treatment for it is placing a stent in the narrowed segment to restore the blood flow. A stent has a mesh-like structure which is crimped onto a cylindrical balloon catheter. This balloon expands with the stent to open the narrowed segment, and then the balloon is deflated and removed, leaving the stent in place to keep the artery open. This stent implantation is very common process such that over 2 million stents are implanted in the United States each year as reported by iData Research \citep{idata2024}. Despite the common usage of stents and the routine nature of the procedure, the mechanical behavior of a single stent after the implantation is still hard to predict. This behavior depends on not just the the design of the stent itself, bit also the arterial wall that is layered, anisotropic, nonlinear, and already damaged by the plaque. It can only be understood by simulating the stent implantation process virtually, and that is what this project is about, making such a simulation applicable and accessible by developing a semi-automated pipeline that can be used for any stent geometry and patient-specific artery model.


\subsection{Problem statement}
\label{sec:problem}
 
A coronary stent has a very complex structure, since it is usually laser-cut out of a thin metal tube. The cut leaves a repeating pattern of ring-shaped units, and each ring consists of slender struts in a zig-zag arrangement that can unfold when the stent is expanded, which is known as the strut pattern. Where the struts are curved is called the crown, and the struts are connected to each other at the crowns to form a ring. The rings are then connected to each other by short links called connectors, so that the stent can stay flexible enough to follow the curvature of the artery, as shown in \cref{fig:stent-geometry} \citep{lin2023}. Therefore, a single stent might contain hundreds of these slender members, and they are all connected to each other, which makes the stent a very complex structure to be modeled at the first place.

\begin{figure}[htb]
  \centering
  \includegraphics[width=0.45\textwidth]{stent_geometry_terminology.png}
  \caption{Terminology used to describe a stent's design \citep{lin2023}.}
  \label{fig:stent-geometry}
\end{figure}


So, this complex and slender structure of the stent is the main reason why the simulation is expensive and hard to be applicable. If the stent is modeled with a common methods for structural analysis, the finite element method with three-dimensional solid elements, the resulting model needs very small elements on every struts along the whole length of the stent in order to capture the characteristics of the struts and crowns correctly. A complete stent model therefore ends up with a fine mesh and a large number of degrees of freedom, and the artery is not even included at this point. The analysis is also strongly nonlinear, since the struts undergo large deformations and large rotations, the stent material deforms elastoplastically, and the struts press against the arterial wall while neighbouring struts can touch each other as well, which is called self-contact. Such a problem is solved incrementally in many small load steps, and even before that, the geometry has to be prepared by generating a conformal volume mesh from a surface of the stent and artery, which is mostly manual work. Therefore, all of these factors make the simulation expensive and time-consuming. Moreover, this challenge is dominated by the stent itself and not by the artery. Because of that, reducing the cost of the simulation by reducing the complexity of the stent structure is one of the main objectives of this project.


\subsection{State of the art}
\label{sec:state-of-the-art}

Finite element modelling for stent deployment has been a topic of research for many years. Various studies have focused on different aspects of stent deployment, including the stresses that different stent designs induce in the arterial wall \citep{lally2005}, or the influence of the balloon model that is used to expand the stent \citep{schiavone2015}. However in order to analyze this mechanical behavior, these studies have always used the classical approach of modelling the stent with three-dimensional solid elements, and the artery is modelled as a hyperelastic solid as well. Even though these models are very successful in capturing the mechanics of the interactions in detail, they are very expensive and time-consuming. In order to reduce this computation time considerably, Hall and Kasper \citep{hall2006} used the fact that the stent is a slender structure, so it can be modelled with 1D beam elements instead of 3D solid ones. This approach is very successful in reducing the computation time, and it still gives acceptable results for the global stress state of the stent after expansion. However, this approach could be just valid if the simplified beam model of the stent is be able to interact with a three-dimensional solid model of the artery. Steinbrecher et al. solved this with a mortar-type formulation that couples 1D beams and 3D solids on non-matching meshes \citep{steinbrecher2020}. This is what makes a mixed-dimensional model practical.

After this successful development on the simplified version of the stent model with mixed-dimensional coupling, Datz. et al. applied this approach to a patient-specific coronary angioplasty \citep{datz2025}. They were be able to simulate the intervention in realistic lesion geometries and found high wall stresses in the severely stenosed areas and at the stent boundaries. Moreover, Ranno et al. used this same approach by also coupling it with hemodynamics, drug release, and tissue growth models to predict the long-term outcome of stent implantation such as restenosis \citep{ranno2025}. However, their stent is a manually built one-ring prototype inside an idealised cylindrical artery. This is sufficient for the verification of the framework, but it is not enough for a real device.

While these studies are very successful in reducing the computation time and making the simulation more applicable, the geometry side of the problem has remained mostly unsolved. Before a simulation can be run, the centerline of the struts has to be extracted from the surface of the stent in order to build a beam model. This is a very challenging task, and it is mostly done manually, which is time-consuming and error-prone. In computer graphics, curve-skeleton extraction is a well known problem where the skeleton of a 2D/3D shape is computed as a graph of curves that run through the centre of its parts. Sharf et al. \citep{sharf2007} extract such a skeleton from point clouds and polygonal meshes, and it stays robust for noisy input and complex topology. Mortier et al. \citep{mortier2008} applied this idea to stents. They generated a beam mesh in a semi-automated way from a triangulated surface, which was obtained from micro-CT images of a real stent sample. However, these studies and algorithms are old and they are not connected to any simulation framework. Therefore, the state of the art is strong on the simulation side, yet what is missing is the step in front of it, namely a reproducible and largely automatic way to turn the surface geometry of a real stent into a simulation-ready one-dimensional model, together with a check that this model actually fits the artery it is going to be deployed in.


\subsection{Objectives and scope}
\label{sec:objectives}

The first aim of this project is to develop an algorithm to turn the surface mesh of a real stent into 1D beam model without any manual CAD editing. In order to makes this algoritm applicable to any stent geometry, some of the features of the stent such as the number of rings, the strut thickness and the diameter are extracted directly from the STL file. The second aim is to prepare a configurated simulation input file for the solver, which is 4C, an open-source multiphysics finite element framework \citep{4c2026}, and how both of these are done is described in \cref{sec:methods}. However, the scope of this project is limited to the centerline extraction pipeline and to the results of only some of the simulation types that the pipeline can build. The two that are reported here are the stent driven alone by prescribed displacements, and the stent opened by contact with an inflating balloon, which follows the paper of Datz et al. \citep{datz2025}. The stent tied into a generated test artery is also built by the pipeline, but it is not reported here. All of these simulations are done with the elastic material for the stent instead of an elastoplastic one, so the results are not directly comparable to the real stent behavior, and this together with the other simplifications made because of the time limitation of this one semester seminar project is discussed in \cref{sec:limitations_and_future}.
